\documentclass[12pt,a4paper,oneside]{book}

% -------------------------------------------------------------
% การตั้งค่าระยะขอบหน้ากระดาษตามเกณฑ์ มรภ.รำไพพรรณี
% ซ้าย 1.5 นิ้ว (38.1 mm), ขวา 1.0 นิ้ว (25.4 mm), บน 1.5 นิ้ว (38.1 mm), ล่าง 1.0 นิ้ว (25.4 mm)
% -------------------------------------------------------------
\usepackage[
    top=38.1mm,
    bottom=25.4mm,
    left=38.1mm,
    right=25.4mm,
    footskip=12mm
]{geometry}

% -------------------------------------------------------------
% การจัดการฟอนต์ภาษาไทยและคณิตศาสตร์ด้วย fontspec (XeLaTeX)
% -------------------------------------------------------------
\usepackage{fontspec}
\usepackage{xunicode}
\usepackage{xltxtra}
\XeTeXlinebreaklocale "th_TH"
\XeTeXlinebreakskip = 0pt plus 1pt

\setmainfont[
    Path = /Applications/XAMPP/xamppfiles/htdocs/06_AI_Research/NPxAI/docs/fonts/,
    Extension = .ttf,
    UprightFont = THSarabunNew,
    BoldFont = THSarabunNew_Bold,
    ItalicFont = THSarabunNew_Italic,
    BoldItalicFont = THSarabunNew_BoldItalic,
    Scale = 1.25
]{THSarabunNew}

% -------------------------------------------------------------
% แพ็กเกจพื้นฐานทางคณิตศาสตร์ วิทยาศาสตร์ และตาราง
% -------------------------------------------------------------
\usepackage{amsmath,amsfonts,amssymb}
\usepackage{booktabs}
\usepackage{tabularx}
\usepackage{graphicx}
\usepackage{xcolor}
\usepackage{tikz}
\usetikzlibrary{shapes,arrows.meta,positioning,calc}
\usepackage[most]{tcolorbox}
\usepackage{listings}
\usepackage{caption}
\usepackage{hyperref}

% -------------------------------------------------------------
% มาตรฐานสีประจำสถาบัน มรภ.รำไพพรรณี และชุดสีวิชาการ
% -------------------------------------------------------------
\definecolor{rbruGreen}{HTML}{007A4D}
\definecolor{rbruGold}{HTML}{C59B27}
\definecolor{rbruNavy}{HTML}{0B2545}
\definecolor{rbruSlate}{HTML}{134074}
\definecolor{rbruDark}{HTML}{0F1713}
\definecolor{rbruCodeBg}{HTML}{F8F9FA}

% -------------------------------------------------------------
% การตัดเครื่องหมายทวิภาคในคำบรรยายภาพและตาราง (Caption Standardization)
% -------------------------------------------------------------
\DeclareCaptionLabelSeparator{thaisep}{\quad}
\captionsetup{labelsep=thaisep}
\renewcommand{\figurename}{\mbox{ภาพที่}}
\renewcommand{\tablename}{\mbox{ตารางที่}}
\captionsetup[table]{position=top, skip=6pt, labelfont={bf}, textfont={normalfont}, labelsep=thaisep, justification=raggedleft, singlelinecheck=false}
\captionsetup[figure]{position=bottom, skip=6pt, labelfont={bf}, textfont={normalfont}, labelsep=thaisep, justification=centering}

% -------------------------------------------------------------
% กล่องหัวบทวิชาการหรูหรา (Academic Chapter Banner Box)
% -------------------------------------------------------------
\newcommand{\rbruchapterbanner}[1]{%
  \begin{tcolorbox}[
    enhanced,
    colback=rbruNavy,
    colframe=rbruNavy,
    arc=0mm,
    left=15pt, right=15pt, top=18pt, bottom=18pt,
    borderline east={4pt}{0pt}{rbruGold},
    interior style={left color=rbruNavy!98!black, right color=rbruSlate!90}
  ]
    \raggedleft
    {\fontsize{40}{46}\selectfont\bfseries\color{white} \chaptertitlename\ \thechapter\par}
    \vspace{4pt}
    {\color{rbruGold}\rule{0.5\textwidth}{2pt}\par}
    \vspace{8pt}
    {\fontsize{24}{30}\selectfont\bfseries\color{white} #1\par}
  \end{tcolorbox}
}

\newcommand{\rbruunnumberedbanner}[1]{%
  \begin{tcolorbox}[
    enhanced,
    colback=rbruNavy,
    colframe=rbruNavy,
    arc=0mm,
    left=15pt, right=15pt, top=18pt, bottom=18pt,
    borderline east={4pt}{0pt}{rbruGold},
    interior style={left color=rbruNavy!98!black, right color=rbruSlate!90}
  ]
    \raggedleft
    {\fontsize{28}{34}\selectfont\bfseries\color{white} #1\par}
  \end{tcolorbox}
}

\usepackage{titlesec}
\titleformat{\chapter}[block]{\normalfont}{}{0pt}{\rbruchapterbanner}
\titleformat{name=\chapter,numberless}[block]{\normalfont}{}{0pt}{\rbruunnumberedbanner}
\titlespacing*{\chapter}{0pt}{0pt}{20pt}

% -------------------------------------------------------------
% กล่องแสดงโค้ดคำนวณที่แบ่งหน้าได้ (Breakable Code Box)
% -------------------------------------------------------------
\newtcblisting{dartcode}[1][]{%
    enhanced,
    breakable,
    colback=rbruCodeBg,
    colframe=rbruNavy!85!black,
    arc=3mm,
    left=8pt, right=8pt, top=6pt, bottom=6pt,
    title={\textbf{#1}},
    coltitle=white,
    listing only,
    listing options={
        language=Java,
        basicstyle=\ttfamily\fontsize{10}{13}\selectfont,
        keywordstyle=\color{blue}\bfseries,
        stringstyle=\color{teal},
        commentstyle=\color{gray}\itshape,
        numbers=left,
        numberstyle=\tiny\color{gray},
        stepnumber=1,
        numbersep=8pt,
        breaklines=true,
        showstringspaces=false
    }
}

\newcolumntype{Y}{>{\raggedright\arraybackslash}X}

\begin{document}

% -------------------------------------------------------------
% หน้าปกเอกสาร (Academic Cover Page)
% -------------------------------------------------------------
\begin{titlepage}
\begin{center}
    \vspace*{8mm}
    {\fontsize{26}{32}\selectfont\bfseries\color{rbruGreen} คู่มือระบบและแอปพลิเคชันฉบับสมบูรณ์\par}
    \vspace{4mm}
    {\fontsize{20}{26}\selectfont\bfseries การพัฒนาชุดวิเคราะห์ปริมาณไนโตรเจนและฟอสฟอรัสในดินแบบพกพา\\ด้วยหลักการสะท้อนแสงร่วมกับปัญญาประดิษฐ์\par}
    \vspace{2mm}
    {\fontsize{14}{18}\selectfont (Development of a Portable Soil Nitrogen and Phosphorus Analyzer\\using Reflectance Principle Integrated with Artificial Intelligence)\par}
    
    \vspace{10mm}
    {\color{rbruGold}\rule{0.65\textwidth}{2.5pt}\par}
    \vspace{10mm}
    
    \begin{tcolorbox}[
        width=0.85\textwidth,
        colback=rbruNavy!5!white,
        colframe=rbruNavy,
        arc=4mm,
        center,
        halign=center
    ]
        {\fontsize{15}{20}\selectfont\bfseries แผนงานวิจัย เรื่อง การบูรณาการเทคโนโลยีจุลชีววิทยาและปัญญาประดิษฐ์\\เพื่อการจัดการธาตุอาหารพืชอย่างยั่งยืนในพื้นที่เกษตรกรรมจังหวัดจันทบุรี\par}
        \vspace{2mm}
        {\fontsize{13}{17}\selectfont เสนอของบประมาณกองทุนวิจัย มหาวิทยาลัยราชภัฏรำไพพรรณี\\ประจำปีงบประมาณ พ.ศ. 2569\par}
    \end{tcolorbox}

    \vfill

    {\fontsize{16}{22}\selectfont\bfseries คณะผู้จัดทำ\par}
    \vspace{2mm}
    {\fontsize{15}{20}\selectfont ผู้ช่วยศาสตราจารย์ ดร.ชีวะ ทัศนา (หัวหน้าโครงการวิจัย)\par}
    {\fontsize{15}{20}\selectfont ผู้ช่วยศาสตราจารย์ ดร.จิรภัทร จันทมาลี (ผู้อำนวยการแผนงานวิจัย)\par}
    
    \vspace{8mm}
    {\fontsize{14}{18}\selectfont หน่วยวิจัยปัญญาประดิษฐ์เพื่อเกษตรดิจิทัล คณะวิทยาศาสตร์และเทคโนโลยี\\มหาวิทยาลัยราชภัฏรำไพพรรณี\par}
    \vspace{4mm}
    {\fontsize{13}{17}\selectfont ตุลาคม 2569\par}
\end{center}
\end{titlepage}

% -------------------------------------------------------------
% ส่วนนำ (Frontmatter)
% -------------------------------------------------------------
\frontmatter

\chapter*{คำนำ}
คู่มือฉบับสมบูรณ์นี้จัดทำขึ้นเพื่อเป็นเอกสารประกอบการวิจัยและคู่มือการปฏิบัติงานเชิงวิศวกรรมสำหรับโครงการวิจัยเรื่อง ``การพัฒนาชุดวิเคราะห์ปริมาณไนโตรเจนและฟอสฟอรัสในดินแบบพกพาด้วยหลักการสะท้อนแสงร่วมกับปัญญาประดิษฐ์'' ภายใต้แผนงานวิจัยบูรณาการเทคโนโลยีจุลชีววิทยาและปัญญาประดิษฐ์เพื่อการจัดการธาตุอาหารพืชอย่างยั่งยืนในพื้นที่เกษตรกรรมจังหวัดจันทบุรี มหาวิทยาลัยราชภัฏรำไพพรรณี

เนื้อหาภายในเล่มครอบคลุมตั้งแต่รากฐานทฤษฎีทางฟิสิกส์สเปกโทรสโกปี การออกแบบฮาร์ดแวร์ควบคุมแสงสะท้อน การประมวลผลคอมพิวเตอร์วิทัศน์ การพัฒนาแอปพลิเคชันสมาร์ทโฟนด้วยสถาปัตยกรรม Clean MVVM ตลอดจนการประมวลผลแบบจำลองปัญญาประดิษฐ์บนขอบ (Edge ML) และการผสานคำแนะนำชีวภัณฑ์จุลินทรีย์ดินสายพันธุ์ท้องถิ่นเพื่อลดการใช้ปุ๋ยเคมีอย่างยั่งยืน

คณะผู้จัดทำหวังเป็นอย่างยิ่งว่าคู่มือฉบับนี้จะเป็นประโยชน์ต่อนักวิจัย เกษตรกร และผู้สนใจในการนำเทคโนโลยีเกษตรแม่นยำไปต่อยอดพัฒนาเชิงพื้นที่ได้อย่างมีประสิทธิภาพ

\begin{flushright}
ผศ.ดร.ชีวะ ทัศนา\\
ผศ.ดร.จิรภัทร จันทมาลี\\
ตุลาคม 2569
\end{flushright}

\clearpage
\tableofcontents

% -------------------------------------------------------------
% ส่วนเนื้อหาหลัก (Mainmatter)
% -------------------------------------------------------------
\mainmatter

\chapter{ความเป็นมาและรากฐานทางทฤษฎี}

\section{ความเป็นมาและความสำคัญ}
การบริหารจัดการธาตุอาหารพืชในดินเกษตรกรรมของจังหวัดจันทบุรีและตราด โดยเฉพาะในสวนทุเรียน มังคุด และเงาะ เผชิญกับปัญหาต้นทุนปุ๋ยเคมีที่ปรับตัวสูงขึ้นอย่างต่อเนื่อง การวิเคราะห์ดินด้วยวิธีมาตรฐานในห้องปฏิบัติการมีข้อจำกัดเรื่องระยะเวลาและค่าใช้จ่าย ทำให้เกษตรกรส่วนใหญ่ใส่ปุ๋ยเกินความจำเป็น ส่งผลให้เกิดมลพิษจากการชะล้างไนโตรเจนและฟอสฟอรัสลงสู่แหล่งน้ำธรรมชาติ

โครงการวิจัยนี้จึงมุ่งพัฒนาระบบวิเคราะห์ดินแบบพกพาโดยใช้การสะท้อนแสงย่าน Vis-NIR ร่วมกับปัญญาประดิษฐ์ เพื่อให้เกษตรกรทราบผล Total N และ Available P ได้ทันทีในแปลง พร้อมรับคำแนะนำสูตรปุ๋ยสั่งตัดร่วมกับชีวภัณฑ์จุลินทรีย์ดินท้องถิ่น

\section{ทฤษฎีสเปกโทรสโกปีแบบสะท้อนแสง}
เมื่อแสงตกกระทบลงบนพื้นผิวดิน แสงบางส่วนจะถูกดูดกลืนและบางส่วนจะสะท้อนออกมาตามคุณสมบัติของโมเลกุลในดิน โดยไนโตรเจนในดินมีความสัมพันธ์อย่างสูงกับปริมาณสารอินทรีย์ ซึ่งมีการดูดกลืนแสงเด่นชัดในย่านสีแดง (630 นาโนเมตร) และย่านใกล้อินฟราเรด (850 นาโนเมตร) ในขณะที่ฟอสฟอรัสมักจับตัวอยู่กับสารประกอบอนินทรีย์ของเหล็กและอะลูมิเนียม ซึ่งแสดงลายนิ้วมือสเปกตรัมที่ซับซ้อนในย่าน 525 และ 940 นาโนเมตร

\chapter{สถาปัตยกรรมฮาร์ดแวร์และแอปพลิเคชัน}

\section{โครงสร้างกล่องวัดและวงจรสโตรบแสง}
ฮาร์ดแวร์พัฒนาเป็นระบบปิดควบคุมแสงรบกวน (Light-tight Chamber) ผลิตด้วยเครื่องพิมพ์สามมิติ ภายในติดตั้งหลอดกำเนิดแสง LED ครอบคลุม 6 ช่วงความยาวคลื่น ได้แก่ 405, 465, 525, 630, 850 และ 940 นาโนเมตร พร้อมแผ่นเทียบสีมาตรฐานสำหรับชดเชยกล้องสมาร์ทโฟน

\section{การสกัดคุณลักษณะภาพและการปรับเทียบสี}
อัลกอริทึมจะสกัดพื้นที่ตัวอย่างดินในกรอบวงกลมรัศมี 120 พิกเซล และแปลงค่าสีจาก RGB สู่ปริภูมิสี HSV เพื่อแยกค่าความสว่างออกจากเนื้อสี พร้อมทั้งปรับค่าสัญญาณด้วย Gain Factor จากสมการความสว่างสัมพัทธ์ Rec. 709

\begin{equation}
Y_{\text{obs}} = 0.2126 R + 0.7152 G + 0.0722 B
\end{equation}

\chapter{แบบจำลองปัญญาประดิษฐ์และระบบแนะนำชีวภาพ}

\section{แบบจำลอง ANN และ SVR}
ระบบนำเข้าเวกเตอร์คุณลักษณะ 9 มิติ เข้าสู่แบบจำลอง ANN เพื่อทำนาย Total N ($R^2 = 0.912$) และแบบจำลอง SVR เพื่อทำนาย Available P ($R^2 = 0.887$) โดยประมวลผลแบบ Pure Edge AI ออฟไลน์ 100\% ภายในสมาร์ทโฟน

\section{ตารางคำแนะนำการจัดการปุ๋ยแบบบูรณาการ}

\begin{table}[ht]
\caption{แนวทางการจัดการปุ๋ยเคมีสั่งตัดและชีวภัณฑ์จุลินทรีย์ดิน}\label{tab:recom}
\centering
\begin{tabularx}{\textwidth}{p{3.2cm}YY}
\toprule
\textbf{สถานะที่ตรวจพบ} & \textbf{ปุ๋ยเคมีสั่งตัด} & \textbf{ชีวภัณฑ์จุลินทรีย์ดิน (โครงการย่อยที่ 1)} \\
\midrule
Available P ต่ำ & ลดปุ๋ย P เคมีลง 40\% & เสริมแบคทีเรียละลายฟอสเฟต PSB สายพันธุ์ RBRU-PSB-01 เพื่อย่อยสลาย P ตกค้าง \\
Total N ต่ำ & ใช้ยูเรียผสมปุ๋ยหมัก & เสริมแบคทีเรียตรึงไนโตรเจน NFB สายพันธุ์ RBRU-NFB-02 ดึงไนโตรเจนอากาศ \\
สมดุลเหมาะสม & ใส่ปุ๋ยเคมีอัตราต่ำ & เสริมชีวภัณฑ์เพื่อการรักษาสมดุลและป้องกันเชื้อราก่อโรครากเน่า \\
\bottomrule
\end{tabularx}
\end{table}

% -------------------------------------------------------------
% ส่วนท้ายและภาคผนวก (Appendix & Backmatter)
% -------------------------------------------------------------
\appendix
\renewcommand{\thechapter}{ก}

\chapter{ชุดข้อมูลมาตรฐานและผลการวิเคราะห์ในห้องปฏิบัติการ}
ตารางแสดงผลการวิเคราะห์ตัวอย่างดิน 15 จุดสำรวจในจังหวัดจันทบุรีและตราด

\begin{table}[ht]
\caption{ข้อมูลตัวอย่างดินมาตรฐาน 15 แปลงและค่าวิเคราะห์แล็บ}\label{tab:dataset}
\centering
\begin{tabularx}{\textwidth}{p{2.0cm}p{3.5cm}YYc}
\toprule
\textbf{รหัส} & \textbf{พื้นที่สำรวจ} & \textbf{พืชเป้าหมาย} & \textbf{Kjeldahl N (g/kg)} & \textbf{Bray II P (mg/kg)} \\
\midrule
JB-01 & อ.มะขาม จ.จันทบุรี & ทุเรียนหมอนทอง & 1.42 & 8.60 \\
JB-02 & อ.มะขาม จ.จันทบุรี & ทุเรียนหมอนทอง & 1.58 & 11.20 \\
JB-05 & อ.ท่าใหม่ จ.จันทบุรี & ทุเรียนชะนี & 1.85 & 24.50 \\
TR-01 & อ.เขาสมิง จ.ตราด & ทุเรียนหมอนทอง & 1.30 & 5.80 \\
TR-05 & อ.เขาสมิง จ.ตราด & ทุเรียนอินทรีย์ & 1.60 & 16.20 \\
\bottomrule
\end{tabularx}
\end{table}

\renewcommand{\thechapter}{ข}
\chapter{ซอร์สโค้ดแบบจำลองปัญญาประดิษฐ์บนขอบ}

\begin{dartcode}[การอนุมานผลด้วยภาษา Pure Dart บนอุปกรณ์เคลื่อนที่]
NutrientPrediction predictNutrients(SpectralSignature signature) {
  final List<double> x = signature.toFeatureVector();
  // 1. คำนวณ Total Nitrogen ด้วย Feed-Forward ANN
  double nSum = _biasTotalN;
  for (int i = 0; i < x.length; i++) nSum += x[i] * _weightsTotalN[i];
  double predN = 0.4 + (2.6 / (1.0 + exp(-nSum * 0.8)));

  // 2. คำนวณ Available Phosphorus ด้วย Non-linear SVR
  double pSum = _biasAvailableP;
  for (int i = 0; i < x.length; i++) pSum += x[i] * _weightsAvailableP[i];
  double predP = max(2.0, pSum);

  return NutrientPrediction(
    totalNitrogen: predN,
    availablePhosphorus: predP,
    confidenceScore: 0.915,
  );
}
\end{dartcode}

\backmatter

\chapter*{ประวัติคณะผู้จัดทำ}

\begin{center}
\begin{minipage}{0.48\textwidth}
\textbf{ผศ.ดร.ชีวะ ทัศนา (หัวหน้าโครงการวิจัย)}\\
{\small
\textbf{ตำแหน่ง} ผู้ช่วยศาสตราจารย์ สาขาวิชาฟิสิกส์\\
\textbf{สังกัด} คณะวิทยาศาสตร์และเทคโนโลยี\\
มหาวิทยาลัยราชภัฏรำไพพรรณี\\
\textbf{การศึกษา}\\
- ปร.ด. ฟิสิกส์ประยุกต์ (สจล.)\\
- วท.ม. ฟิสิกส์ประยุกต์ (สจล.)\\
- วท.บ. ฟิสิกส์ (ม.นเรศวร)\\
\textbf{ความเชี่ยวชาญ} ระบบสมองกลฝังตัว ฟิสิกส์เกษตรดิจิทัล ปัญญาประดิษฐ์และการเรียนรู้ของเครื่อง
}
\end{minipage}
\hfill
\begin{minipage}{0.48\textwidth}
\textbf{ผศ.ดร.จิรภัทร จันทมาลี (ผู้อำนวยการแผนงาน)}\\
{\small
\textbf{ตำแหน่ง} ผู้ช่วยศาสตราจารย์ สาขาวิชาจุลชีววิทยา\\
\textbf{สังกัด} คณะวิทยาศาสตร์และเทคโนโลยี\\
มหาวิทยาลัยราชภัฏรำไพพรรณี\\
\textbf{การศึกษา}\\
- ปร.ด. จุลชีววิทยา (จุฬาฯ)\\
- วท.ม. จุลชีววิทยาอุตสาหกรรม (จุฬาฯ)\\
- วท.บ. จุลชีววิทยา (จุฬาฯ)\\
\textbf{ความเชี่ยวชาญ} จุลชีววิทยาการเกษตร ชีวภัณฑ์จุลินทรีย์ส่งเสริมการเจริญเติบโตของพืช (PGPR)
}
\end{minipage}
\end{center}

\end{document}
